\section{Signed perturbations on the retained head}
\label{sec:signed-head-lower}

The retained \(\ell_{2/3}\) term, deletion of a largest atom, and
clipped \(2/3\)-power shifts appear in Valiant and Valiant's
instance-optimal testing analysis \cite{valiantvaliant2017}.
Chhor and Carpentier formulate the corresponding fixed-sample
separation radius and use signed priors in their bulk lower bound
\cite[Corollary~1 and Section~5.1]{chhorcarpentier2022}. We use
capped, paired shifts and keep the original head probabilities. This
preserves the total head mass and lets the same prior be composed with
the independent tail construction.

List the positive head probabilities as \(q_1\ge\cdots\ge q_d>0\).
They need not sum to one. Set \(w=1-\sum_{i=1}^{d}q_i\) and
\(\Delta=e-w\). On the reduced alphabet, write \(q^\sharp\) for the
law with these \(d\) probabilities and an additional tail marker of
mass \(w\). For \(n\ge1\), use the head profile
\begin{equation}
 P(n)=\min_{1\le j\le d}
 \left\{
  \frac{\left(\sum_{i=2}^{j}q_i^{2/3}\right)^{3/4}}{\sqrt n}
  +\sum_{i=j+1}^{d}q_i
 \right\}.
 \label{eq:signed-head-profile}
\end{equation}
If the normalized head law is \(h_i=q_i/A\), where \(A=1-w>0\),
then this profile equals \(A P_h(nA)\) with \(P_h\) given by the
same formula on \(h\). Both terms scale by \(A\) under this
substitution. We will use \(P(n)\) directly.

\paragraph{A capped shift with a quartic budget.}
For \(i=2,\ldots,d\), choose
\[
 z_i=\min(q_i,\lambda q_i^{2/3}),
 \qquad
 \sum_{i=2}^{d}\frac{z_i^4}{q_i^2}\le\frac1{n^2}.
\]
If \(\sum_{i=2}^{d}q_i^2\le n^{-2}\), take \(\lambda\) large
enough that every \(z_i=q_i\). Otherwise the budget is attained with
equality for some \(\lambda>0\), by continuity between
\(\lambda=0\) and \(\lambda=\max_{i\ge2}q_i^{1/3}\). In the
singleton case all sums are empty and the saturated choice applies.
We extend the same definition of \(z_i\) to \(i=1\) for ordering
purposes and write \(S=\sum_{i=2}^{d}z_i\).

For any cutoff \(j\), H\"older's inequality and the quartic budget
give
\[
 \sum_{i=2}^{j} z_i
 \le \left(\sum_{i=2}^{j}z_i^4/q_i^2\right)^{1/4}
       \left(\sum_{i=2}^{j}q_i^{2/3}\right)^{3/4}
 \le n^{-1/2}\left(\sum_{i=2}^{j}q_i^{2/3}\right)^{3/4}.
\]
The remaining shift mass is at most the remaining reference mass.
Thus \(S\le P(n)\). In the all-saturated case the cutoff \(j=1\)
also gives \(P(n)\le S\). Otherwise put
\(I=\{i\ge2:q_i>\lambda^3\}\) and
\(I^c=\{2,\ldots,d\}\setminus I\). The retained indices \(I\)
form a prefix. Write
\[
 W=\sum_{i\in I}q_i^{2/3},\qquad
 U=\sum_{i\in I^c}q_i,\qquad R=\sum_{i\in I^c}q_i^2.
\]
Here \(S=\lambda W+U\), \(n^{-2}=\lambda^4W+R\), and
\(R\le\lambda^3 U\). The fourth power of the bulk term at this
cutoff satisfies
\[
 \frac{W^3}{n^2}
 \le (\lambda W)^4+(\lambda W)^3U
 \le (\lambda W+U)^4.
\]
Consequently \(S\le P(n)\le2S\). Atoms tied at
\(\lambda^3\) may stay outside the prefix.

\paragraph{The paired prior.}
Use the disjoint pairs \((1,2),(3,4),\ldots\). For each pair \(j\)
draw an independent uniform sign \(\sigma_j\in\{-1,1\}\) and add
\(\sigma_j b_j\) to coordinate \(2j-1\) and subtract it from
coordinate \(2j\), where \(b_j=z_{2j}/8\). Write \(p_\sigma^\sharp\)
for the resulting reduced law. Since \(b_j\le q_{2j}\le q_{2j-1}\),
the shifts give a nonnegative normalized law. They leave the marker
and any unpaired label unchanged. In particular the construction
permits a marker of mass zero. The function
\(t\mapsto\min(t,\lambda t^{2/3})\) is nondecreasing, so the \(z_i\)
are sorted. Each odd index \(i\ge3\), including a possible unpaired
last index, has \(z_i\le z_{i-1}\). For every sign vector the
unconditional head deficit is therefore
\begin{equation}
 L_H(p_\sigma^\sharp,q^\sharp)=\sum_j b_j
   \ge S/16\ge P(n)/32 .
 \label{eq:signed-head-deficit}
\end{equation}

Let \(M=\mathbb E_\sigma (p_\sigma^\sharp)^{\otimes n}\) and
\(Q=(q^\sharp)^{\otimes n}\) on ordered \(n\)-tuples. The
perturbations have no mass outside the null support, including at a
zero marker. With independent sign vectors \(\sigma,\tau\), product
factorization of the likelihood inner product gives
\begin{align*}
 1+\chi^2(M,Q)
 &=\mathbb E_{\sigma,\tau}
   \left(\sum_{a:q^\sharp_a>0}
      \frac{(p_\sigma^\sharp)_a(p_\tau^\sharp)_a}{q^\sharp_a}
   \right)^n  \\
 &=\mathbb E_{\sigma,\tau}
   \left(1+\sum_j c_j\sigma_j\tau_j\right)^n,
 & c_j&=b_j^2\left(\frac1{q_{2j-1}}+\frac1{q_{2j}}\right).
\end{align*}
The linear shift terms cancel within each pair, and the unchanged
labels contribute their null masses to the constant term. The
quantity raised to the power \(n\) is nonnegative. The inequality
\(1+u\le e^u\) and independence of the sign products now give
\[
 1+\chi^2(M,Q)
 \le \prod_j\cosh(nc_j)
 \le \exp\!\left(\frac{n^2}{2}\sum_j c_j^2\right)
 \le \exp(1/2048).
\]
The middle step follows from the power series and
\((2k)!\ge2^k k!\). The last follows from
\(c_j\le 2b_j^2/q_{2j}\) and the quartic budget. Thus
\(\operatorname{TV}(M,Q)\le\frac12\sqrt{e^{1/2048}-1}<1/8\)
by Cauchy--Schwarz. For the numerical inequality one can use
\(e^t\le(1-t)^{-1}\) at \(t=1/2048\).

\paragraph{Consequences at the exact radius.}
Suppose the tail construction gives a prior of tail laws of mass
\(w\) whose reference overlap is at most \(x\) for every draw of the
prior. Suppose also that when the head probabilities are unchanged,
its ordered \(n\)-sample mixture is within \(1/8\) of the full
null. When \(x\ge w\), the unchanged tail is such a prior. In
particular the argument covers \(w=0\), where the marker is never
observed. For positive \(w\), choose the tail prior independently
of the head signs and replace each marker by an independent draw
from the chosen conditional tail law, given that law. This is a
single common randomized channel on ordered \(n\)-tuples. It
contracts total variation, so the full mixture of paired alternatives
is within \(1/8+1/8=1/4\) of the null.

If \(P(n)>32(\Delta+x)\), equation~\eqref{eq:signed-head-deficit}
gives \(L_H>\Delta+x\) for every paired alternative. Marker
replacement leaves this deficit unchanged and gives
\(O_T=\sum_{i\in T}\min(p_i,q_i)\le x\). Hence every lifted law
has distance \(w+L_H-O_T>e\). A test with each of its two errors
at most \(1/3\) would then have a rejection-probability gap of at
least \(1/3\) between the null and this mixture. That is impossible
at mixture TV at most \(1/4\). Thus a successful \(n\)-draw test
necessarily satisfies
\begin{equation}
 P(n)\le32(\Delta+x)
 \label{eq:signed-head-necessity}
\end{equation}
whenever the stated tail-prior bounds hold. The radius in this
comparison is the original strict threshold \(e\).
